\documentclass{article}

\usepackage{amsmath,amssymb}
\usepackage{mathrsfs}
\usepackage{xurl}
\usepackage[hidelinks]{hyperref}
\setlength{\emergencystretch}{2em}

% Shared commands for notes in docs/paper-gaps/.

\DeclareUnicodeCharacter{2080}{\ifmmode{}_{0}\else\textsubscript{0}\fi}
\DeclareUnicodeCharacter{2081}{\ifmmode{}_{1}\else\textsubscript{1}\fi}
\DeclareUnicodeCharacter{2082}{\ifmmode{}_{2}\else\textsubscript{2}\fi}
\DeclareUnicodeCharacter{2099}{\ifmmode{}_{n}\else\textsubscript{n}\fi}

\newcommand{\Lean}{\textsc{Lean}}
\ExplSyntaxOn
\NewDocumentCommand{\leanid}{m}{
  \ifmmode
    \textsf{\detokenize{#1}}
  \else
    \group_begin:
    \urlstyle{sf}
    \tl_set:Nn \l_tmpa_tl {#1}
    \tl_replace_all:Nnn \l_tmpa_tl {\_} {_}
    \exp_args:NV \path \l_tmpa_tl
    \group_end:
  \fi
}
\ExplSyntaxOff
\providecommand{\tr}{\operatorname{tr}}
\newcommand{\tnleanIssueBase}{https://github.com/LionSR/TNLean/issues/}
\newcommand{\tnleanPrBase}{https://github.com/LionSR/TNLean/pull/}
\newcommand{\tnleanDocBase}{https://sirui-lu.com/TNLean/docs/}
\newcommand{\ghissue}[1]{\href{\tnleanIssueBase#1}{GitHub issue \##1}}
\newcommand{\ghpr}[1]{\href{\tnleanPrBase#1}{PR \##1}}
\newcommand{\doclink}[1]{\href{\tnleanDocBase#1.html}{\path{#1}}}

% Dirac notation and the identity operator, as in blueprint/src/macros/common.tex.
% \providecommand keeps note-local definitions authoritative.
\usepackage{dsfont}
\providecommand{\ket}[1]{|#1\rangle}
\providecommand{\bra}[1]{\langle #1|}
\providecommand{\braket}[2]{\langle #1 | #2 \rangle}
\providecommand{\ketbra}[2]{|#1\rangle\!\langle #2|}
\providecommand{\Id}{\mathds{1}}
\providecommand{\one}{\mathds{1}}
\providecommand{\C}{\mathbb{C}}
\providecommand{\MN}[1]{M_{#1}(\C)}
\providecommand{\MD}{M_D(\C)}

% External referencing for paper sources.  The build helper writes sanitized
% auxiliary files under build/paper-aux/ and excludes duplicated source labels.
\usepackage{xr}

% Import a generated paper auxiliary file with a caller-supplied label prefix.
% The candidate paths support builds from docs/paper-gaps/, the repository root,
% and build/paper-aux/ itself.
\newcommand{\importpaperaux}[2]{%
  \IfFileExists{../../build/paper-aux/#2.aux}{%
    \externaldocument[#1]{../../build/paper-aux/#2}%
  }{%
    \IfFileExists{build/paper-aux/#2.aux}{%
      \externaldocument[#1]{build/paper-aux/#2}%
    }{%
      \IfFileExists{#2.aux}{%
        \externaldocument[#1]{#2}%
      }{}%
    }%
  }%
}

% General external reference macros
\newcommand{\paperref}[2]{\ref{#1-#2}}
\newcommand{\papereqref}[2]{(\ref{#1-#2})}

% CPSV16 source (1606.00608 / MPDO-22-12-17-2.tex) external document and shortcuts
\importpaperaux{cpsv16-}{1606.00608/MPDO-22-12-17-2-tnlean}
\newcommand{\cpsvref}[1]{\paperref{cpsv16}{#1}}
\newcommand{\cpsveqref}[1]{\papereqref{cpsv16}{#1}}

% Verdict marker: \gapnote{<kind>}{<status>}. Kind is the policy
% classification (clarification, local-correction, scope-restriction,
% unfaithful, false-source, open-gap); status is open, wip (elimination
% underway), resolved, or historical. Machine-read by the site index;
% prints a verdict line.
\newcommand{\gapnote}[2]{\par\noindent\textbf{Verdict:} #1 (#2).\par}

\providecommand{\Fin}[1]{\{0,\ldots,#1{-}1\}}

\title{The Martingale Principal-Angle Estimate for MPS Parent Hamiltonians}
\author{}
\date{2026-06-04}

\begin{document}
\maketitle
\gapnote{scope-restriction}{open}

\section{The source assertion}

The martingale argument of
Cirac--P\'erez-Garc\'ia--Schuch--Verstraete can be compared directly with the
quantitative estimate used in the formal parent-Hamiltonian
proof~\cite[Section~IV.C]{Cirac2021Matrix}.\footnote{The finite-row reduction
    is tracked in
    \href{https://github.com/LionSR/TNLean/issues/1044}{issue \#1044}.
    The blocked MPS ground-space estimate is tracked in
    \href{https://github.com/LionSR/TNLean/issues/5455}{issue \#5455}, and the
    broader parent-Hamiltonian theorem is tracked in
    \href{https://github.com/LionSR/TNLean/issues/460}{issue \#460}.}
The relevant source paragraph is \path{sec:4:hams-gaps}, which contains the
inequalities \path{eq:4:martingale-1} and
\path{eq:4:martingale-2}~\cite[Section~IV.C]{Cirac2021Matrix}.

Let
\begin{align}
    H&=\sum_i h_i
    \label{eq:cpgsv21_martingale_overlap_hamiltonian_sum}
\end{align}
be a frustration-free Hamiltonian whose local terms $h_i$ are orthogonal
projections. At the quadratic-form level, a gap lower bound $\gamma>0$ is
equivalent to
\begin{align}
    H^2&\geq\gamma H.
    \label{eq:cpgsv21_martingale_overlap_gap_quadratic_form}
\end{align}
Because $h_i^2=h_i$, the expansion of $H^2$ contains diagonal terms,
overlapping products, and non-overlapping products. The non-overlapping
products are non-negative. It therefore suffices to prove, for every
overlapping pair, the anticommutator estimate
\begin{align}
    h_i h_j+h_j h_i
    &\geq-c_{ij}(1-\gamma)(h_i+h_j),
    \label{eq:cpgsv21_martingale_overlap_source_anticommutator}
\end{align}
where the coefficients $c_{ij}$ have row sum one in
\path{eq:4:martingale-2} of
Ref.~\cite[Section~IV.C]{Cirac2021Matrix}. The formal row-sum argument permits
the more general condition that every row sum is at most one.

The same source paragraph interprets
(\ref{eq:cpgsv21_martingale_overlap_source_anticommutator}) through the
minimum nonzero angle between ground spaces of overlapping regions. It cites
the martingale method and the MPS parent-Hamiltonian results of
Fannes--Nachtergaele--Werner, Nachtergaele, and Kastoryano--Lucia to conclude
that suitable blocking gives a uniform positive gap
\cite{Fannes1992Finitely,Nachtergaele1996Spectral,Kastoryano2018Martingale}.

To distinguish ground-space and excitation projections, let $q_i$ and $q_j$
project onto $\ker h_i$ and $\ker h_j$, respectively, and let $P$ project onto
their common ground space. For overlap length $\ell$, define
\begin{align}
    \delta(\ell)
    &:=\lVert q_iq_j-P\rVert,
    \label{eq:cpgsv21_martingale_overlap_ground_projector_defect}\\
    \operatorname{ran}P
    &=\ker h_i\cap\ker h_j
      =\ker(h_i+h_j).
    \label{eq:cpgsv21_martingale_overlap_common_ground_space}
\end{align}
The symbols $h_i,h_j$ in the corresponding review display must be read as the
ground-space projectors $q_i,q_j$, not as the excitation projections in
(\ref{eq:cpgsv21_martingale_overlap_source_anticommutator}). Substituting the
excitation projections into
(\ref{eq:cpgsv21_martingale_overlap_ground_projector_defect}) would make the
norm at least one on the common ground space. For the ground-space projectors,
the cited result gives exponential decay of $\delta(\ell)$ with universal
constants whenever the frustration-free Hamiltonian is gapped. This decay
implies the martingale condition because both quantities depend on the
relevant principal angles~\cite{Kastoryano2018Martingale}. The review does not
state the numerical constants or the finite cyclic blocking convention needed
below. In particular,
(\ref{eq:cpgsv21_martingale_overlap_common_ground_space}) identifies the
common ground space but gives no quantitative angle estimate.

For finite regions $A$ and $B$, Kastoryano--Lucia define
\begin{align}
    \delta(A,B)
    &:=
      \lVert(P_A-P_{A\cup B})(P_B-P_{A\cup B})\rVert,
    \label{eq:cpgsv21_martingale_overlap_kl_defect_definition}\\
    \delta(A,B)
    &=\lVert P_AP_B-P_{A\cup B}\rVert
    \label{eq:cpgsv21_martingale_overlap_kl_defect_product}
\end{align}
under frustration-freeness, and identify this quantity with the cosine of the
first principal angle between the reduced ground-space projectors
\cite[Definition~2 and Remark~2]{Kastoryano2018Martingale}. Nachtergaele's
conditions C3 and C3$'$ instead require
\begin{align}
    \lVert G_{\Lambda_{n+1}\setminus\Lambda_{n-l}}E_n\rVert
    &\leq\varepsilon_l,
    &
    \varepsilon_l
    &<\frac{1}{\sqrt{l+1}},
    \label{eq:cpgsv21_martingale_overlap_nachtergaele_c3}\\
    \lVert G_{\Lambda_{n+p}\setminus\Lambda_{n-l}}E_n^{(p)}\rVert
    &\leq\eta_l,
    &
    \eta_l
    &<\frac{1}{\sqrt{2}}.
    \label{eq:cpgsv21_martingale_overlap_nachtergaele_c3_prime}
\end{align}
These are Equations~(2.4)--(2.5) of
Ref.~\cite{Nachtergaele1996Spectral}; Theorem~3 of the same reference gives
the corresponding gap lower bounds. Section~6 verifies C3$'$ for GVBS/MPS
Hamiltonians using projection estimates inherited from
Ref.~\cite{Fannes1992Finitely}. The formal development instead proves a
sufficient range-two estimate after physical blocking. The optimized FNW
numerator and the identification of its source constants remain open.

\subsection{Unit bound at the full-window volume}

Fix an MPS tensor $A$ and an interaction length $L>0$. Let
$P_{\mathcal G_L^{\mathrm{ES}}(A)}$ be the orthogonal projector onto the
length-$L$ MPS ground space, and define the excitation projector by
\begin{align}
    P_L^{\mathrm{ES}}(A)
    &:=\Id-P_{\mathcal G_L^{\mathrm{ES}}(A)}.
    \label{eq:cpgsv21_martingale_overlap_excitation_projector}
\end{align}
For a chain of length $N$ and $i+L\leq N$, let
$h_{i,\mathrm{ES}}(A,L)$ be the copy of this projector on
$[i,i+L-1]$, extended by the identity outside that window. Define
\begin{align}
    I_{L,N}
    &:=
      \{i\in\{0,\ldots,N-1\}:i+L\leq N\},
    \label{eq:cpgsv21_martingale_overlap_compatible_starts}\\
    H^{\mathrm{open}}_{N,L}(A)
    &:=
      \sum_{i\in I_{L,N}}h_{i,\mathrm{ES}}(A,L).
    \label{eq:cpgsv21_martingale_overlap_open_hamiltonian}
\end{align}
At $N=L$, this Hamiltonian contains one orthogonal
projection.\footnote{This is the Blueprint theorem ``Full-window compatible
    interaction,'' labelled
    \path{thm:open_parent_hamiltonian_full_window}; its formal declaration is
    \leanid{MPSTensor.openParentHamiltonianES_self_eq_parentInteractionES}.}
Consequently,
\begin{align}
    H^{\mathrm{open}}_{L,L}(A)
    &=P_L^{\mathrm{ES}}(A).
    \label{eq:cpgsv21_martingale_overlap_full_window_identity}
\end{align}
It preserves the norm on the orthogonal complement of its kernel. This
identity covers only the single-window case and gives no estimate for $N>L$.
The general-volume result also needs the open-chain ground-space intersection,
a quantitative projector-defect estimate, and the finite-row martingale
reduction.

\subsection{Open-chain projector geometry}

The open-chain geometry is independent of the later cyclic specialization.
Inside the Hilbert space on $\Lambda_{n+1}$, the formal subspaces
\leanid{MPSTensor.openChainLeftGroundSpaceES} and
\leanid{MPSTensor.openChainTailGroundSpaceES} are the ranges of $G_{\Lambda_n}$ and
$G_{\Lambda_{n+1}\setminus\Lambda_{n-l}}$. They use the existing last-site
and fixed-prefix restriction maps. Under $L_0$-block injectivity, with
$n=K+L_0$ and $l=L_0$, the theorem
\leanid{MPSTensor.openChainLeft_inf_tailGroundSpaceES} proves
\begin{align}
    U\cap V
    &=\mathcal G_{K+L_0+1}^{\mathrm{ES}}(A).
    \label{eq:cpgsv21_martingale_overlap_open_chain_intersection}
\end{align}

The theorem
\leanid{Submodule.isFriedrichsBound_of_norm_starProjection_comp_sub_inf_starProjection_le}
turns
\begin{align}
    \lVert P_UP_V-P_{U\cap V}\rVert
    &\leq\eta
    \label{eq:cpgsv21_martingale_overlap_friedrichs_projector_bound}
\end{align}
into a Friedrichs bound. Combined with the complementary-projection theorem,
it gives
\leanid{MPSTensor.re_inner_openChain_anticommutator_ge_neg_of_groundProjection_defect}.
Its hypothesis is the C3 defect
\begin{align}
    \lVert
      G_{\Lambda_{n+1}\setminus\Lambda_{n-l}}G_{\Lambda_n}
      -G_{\Lambda_{n+1}}
    \rVert
    &\leq\eta,
    \label{eq:cpgsv21_martingale_overlap_open_chain_c3_defect}
\end{align}
using the identity following condition C3 in Equation~(2.4) of
Ref.~\cite{Nachtergaele1996Spectral}. Equation~(2.5) of the same reference is
the distinct batched condition C3$'$ involving $E_n^{(p)}$.

The numerical open-chain step is complete. The geometric estimate is uniform
in the left length, and
\leanid{MPSTensor.IsPrimitiveMPS.exists_openChain_groundProjection_defect_lt_c3_threshold}
chooses a block-injective $l>1$, measured in sites of the input tensor, and
$\varepsilon_l\geq0$ such that
\begin{align}
    \varepsilon_l
    &<\frac{1}{\sqrt{l+1}}.
    \label{eq:cpgsv21_martingale_overlap_c3_threshold}
\end{align}
The theorem
\leanid{MPSTensor.IsPrimitiveMPS.exists_re_inner_openChain_anticommutator_ge_c3_threshold}
combines Nachtergaele's projector-defect identity with the two-projection
estimate and proves the corresponding anticommutator inequality for every left
length. Exact blocked three-site identities then connect the three-block defect
to the local terms of the blocked tensor. No excitation-compression estimate
is used.

\subsection{Three-site to cyclic transport}

Let $h_0,h_1$ be adjacent two-site terms on a three-site Hilbert space. Assume
\begin{align}
    h_0h_1+h_1h_0
    &\geq-c(h_0+h_1)
    \label{eq:cpgsv21_martingale_overlap_three_site_anticommutator}
\end{align}
in quadratic form. The theorems
\leanid{MPSTensor.re_inner_localTermES_adjacent_twoSite_forward_ge_of_threeSite}
and
\leanid{MPSTensor.re_inner_localTermES_adjacent_twoSite_backward_ge_of_threeSite}
transport
(\ref{eq:cpgsv21_martingale_overlap_three_site_anticommutator}), with the same
coefficient $c$, to every forward and backward adjacent cyclic pair on every
chain of length $N\geq3$. The proof decomposes the configuration basis into a
cyclic three-site window and its ordered complement, applies the three-site
estimate fiberwise through cyclic restriction, and sums over the complement.
Arbitrary window starts cover the periodic cut; the backward theorem starts
the forward window at the cyclic predecessor.

The exact identification
\leanid{MPSTensor.blockTensor_threeSite_openChain_defect_norm_eq} equates the
three-site open-chain projector defect of
$B=\operatorname{block}_p(A)$ with the whole-block defect on three consecutive
$p$-site regions of $A$. Together with
\leanid{MPSTensor.IsPrimitiveMPS.exists_threeBlock_wholeIncrement_defect_le_seven_sixteenths},
it gives
\begin{align}
    c&=\frac{7}{16}
    \label{eq:cpgsv21_martingale_overlap_three_block_coefficient}
\end{align}
for adjacent three-site terms of $B$. The two transport theorems cover all
overlapping off-diagonal cyclic pairs, including pairs across the periodic cut.
This argument neither claims the optimized FNW coefficient nor uses an
excitation-compression surrogate.

\section{Finite-row martingale reduction}

The formal development separates the general finite-row algebra from the
completed range-two blocked MPS case. Fix $L>1$ and an $N$-site periodic chain
with $N\geq2L$. Let $p_i$ be the Hilbert-space representative of the
length-$L$ local parent-Hamiltonian projection on the cyclic window starting at
$i$. Write $i\sim_Lj$ when the two windows intersect and $i\neq j$. Each row
has at most
\begin{align}
    m&=2(L-1)
    \label{eq:cpgsv21_martingale_overlap_overlap_row_cardinality}
\end{align}
such indices. Choose
\begin{align}
    c_{ij}
    &:=
    \begin{cases}
        \dfrac{1}{2(L-1)},& i\sim_Lj,\\
        0,&\text{otherwise}.
    \end{cases}
    \label{eq:cpgsv21_martingale_overlap_uniform_row_coefficient}
\end{align}
Then
\begin{align}
    \sum_jc_{ij}
    &\leq1.
    \label{eq:cpgsv21_martingale_overlap_row_sum_bound}
\end{align}

Set
\begin{align}
    \gamma_L
    &:=\frac{1}{4L}.
    \label{eq:cpgsv21_martingale_overlap_special_gap_constant}
\end{align}
A sufficient ordered-overlap estimate is
\begin{align}
    \operatorname{Re}\langle p_iv,p_jv\rangle
    &\geq
      -\left(1-\frac{1}{4L}\right)
       \frac{1}{2(L-1)}
       \operatorname{Re}\langle p_iv,v\rangle
    \label{eq:cpgsv21_martingale_overlap_ordered_cross_term}
\end{align}
for every $N\geq2L$, every $i\sim_Lj$, and every vector $v$. The formal
row-sum argument then proves
\begin{align}
    H_{N,L}^2
    &\geq\gamma_LH_{N,L},
    \label{eq:cpgsv21_martingale_overlap_finite_row_gap}
\end{align}
and hence gives a norm lower bound with constant $\gamma_L$ on the orthogonal
complement of the ground space.

A stronger sufficient input is the directed norm-compression estimate
\begin{align}
    \lVert p_ip_jv\rVert
    &\leq
      \left(1-\frac{1}{4L}\right)
      \frac{1}{2(L-1)}
      \lVert p_iv\rVert.
    \label{eq:cpgsv21_martingale_overlap_special_norm_compression}
\end{align}
Because $p_i$ is an orthogonal projection,
\begin{align}
    \operatorname{Re}\langle p_iv,p_jv\rangle
    &=\operatorname{Re}\langle p_iv,p_ip_jv\rangle,
    \label{eq:cpgsv21_martingale_overlap_projected_cross_term}\\
    \operatorname{Re}\langle p_iv,p_ip_jv\rangle
    &\geq-\lVert p_iv\rVert\lVert p_ip_jv\rVert,
    \label{eq:cpgsv21_martingale_overlap_cross_term_cauchy_schwarz}\\
    \operatorname{Re}\langle p_iv,v\rangle
    &=\lVert p_iv\rVert^2.
    \label{eq:cpgsv21_martingale_overlap_projection_quadratic_form}
\end{align}
Thus
(\ref{eq:cpgsv21_martingale_overlap_special_norm_compression}) implies
(\ref{eq:cpgsv21_martingale_overlap_ordered_cross_term}). The formal
development also contains the source anticommutator row-sum argument directly,
without this stronger one-sided estimate.\footnote{The formal anchors are
    \leanid{ProjectionGeometry.crossTerm_sum_bound_of_anticommutator_rowCol},
    \leanid{ProjectionGeometry.quadraticForm_sum_projections_of_anticommutator_rowCol}, and
    \leanid{MPSTensor.parentHamiltonianES_gap_bound_of_anticommutator_local_term_bounds}.
    The finite-overlap reductions are
    \leanid{ProjectionGeometry.quadraticForm_sum_projections_of_finite_overlap_anticommutator},
    \leanid{MPSTensor.parentHamiltonianES_gap_bound_of_finite_overlap_anticommutator},
    and
    \leanid{MPSTensor.parentHamiltonianES_gap_bound_of_cyclic_window_overlap_anticommutator}.
    The final gap statement is
    \leanid{MPSTensor.parentHamiltonian_gapped_of_anticommutator} in
    \path{TNLean/MPS/ParentHamiltonian/Martingale/Gap.lean}.
    The cyclic-window compression reductions are
    \leanid{MPSTensor.parentHamiltonianES_gap_bound_of_cyclic_window_ordered_cross_term},
    \leanid{MPSTensor.parentHamiltonianES_gap_bound_of_cyclic_window_overlap_norm_bound},
    and
    \leanid{MPSTensor.parentHamiltonianES_gap_bound_of_cyclic_window_overlap_norm_bound_of_le}
    in \path{TNLean/MPS/ParentHamiltonian/Martingale/Reduction.lean}.
    The constant-$\eta$ reduction is
    \leanid{MPSTensor.parentHamiltonianES_gap_bound_of_cyclic_window_overlap_norm_bound_of_lt}
    in the same file; the public spectral-gap statement is
    \leanid{MPSTensor.parentHamiltonian_gapped_of_overlap_norm_constant} in
    \path{TNLean/MPS/ParentHamiltonian/Martingale/Gap.lean}.
    The projection-geometry results are in
    \path{QICLean/Analysis/ProjectionGeometry.lean}. The declaration names
    distinguish the ordered cross-term estimate from the sufficient
    norm-compression estimates. This note reserves ``anticommutator estimate''
    for the source condition and ``norm compression'' for the stronger local
    condition.}

The finite-overlap cyclic anticommutator formulation is the formal condition
that matches the source. A flexible norm-compression variant assumes
$\eta\geq0$ and
\begin{align}
    2\eta(L-1)
    &<1,
    \label{eq:cpgsv21_martingale_overlap_eta_smallness}\\
    \lVert p_ip_jv\rVert
    &\leq\eta\lVert p_iv\rVert
    \label{eq:cpgsv21_martingale_overlap_eta_norm_compression}
\end{align}
for all overlapping cyclic-window pairs. The row-sum argument then gives
\begin{align}
    \gamma
    &=1-2\eta(L-1)>0.
    \label{eq:cpgsv21_martingale_overlap_eta_gap}
\end{align}
Equation~(\ref{eq:cpgsv21_martingale_overlap_special_norm_compression}) is one
possible specialization. The finite-row comparison only requires a compression
constant below $1/(2(L-1))$ in the chosen blocking convention.

\section{What remains to compare}

The source and formal arguments agree qualitatively: both expand the
frustration-free Hamiltonian, use positivity for non-overlapping terms, assign
row-sum coefficients to overlapping terms, and control those terms through
principal-angle geometry. The unresolved comparison concerns constants and
blocking conventions.

The source-matching formal theorem assumes
\begin{align}
    p_ip_j+p_jp_i
    &\geq
      -\left(1-\frac{1}{4L}\right)
       \frac{1}{2(L-1)}(p_i+p_j)
    \label{eq:cpgsv21_martingale_overlap_cyclic_anticommutator_target}
\end{align}
for every overlapping off-diagonal pair. The review states the corresponding
criterion with general row-sum coefficients and invokes blocking and principal
angles to obtain a gap~\cite[Section~IV.C]{Cirac2021Matrix}. It does not
identify the estimates of
Refs.~\cite{Fannes1992Finitely,Nachtergaele1996Spectral,Kastoryano2018Martingale}
with this coefficient and cyclic-window convention. A proof could instead
derive a different row-sum anticommutator estimate and adjust the final gap
constant.

The formal theorem is uniform for all $N\geq2L$ and uses length-$L$ cyclic
windows in its finite-chain Hilbert-space model. A source proof that changes
the blocking can still establish an existential MPS gap theorem, but it does
not immediately imply
(\ref{eq:cpgsv21_martingale_overlap_cyclic_anticommutator_target}) for every
$L$.

The optimized FNW coefficient and its source constants remain unmatched. This
comparison is unnecessary for the proved range-two blocked theorem: the exact
three-site identification gives $7/16$, and cyclic transport covers every
overlap. The theorem
\leanid{Submodule.re_inner_anticommutator_starProjection_orthogonal_ge_neg_of_friedrichs_bound}
shows that a Friedrichs bound $\eta$ for the two ground-space ranges, after
removing their common intersection, gives
\begin{align}
    h_ih_j+h_jh_i
    &\geq-\eta(h_i+h_j)
    \label{eq:cpgsv21_martingale_overlap_friedrichs_anticommutator}
\end{align}
for the complementary excitation projections. The remaining question is
whether the cited estimates recover the optimized numerator and source
constants in the present Gram and blocking conventions. Existence of a
finite-dimensional principal-angle decomposition alone does not provide the
required uniform numerical bound.

The all-vector condition
(\ref{eq:cpgsv21_martingale_overlap_eta_norm_compression}) remains a stronger
conditional route. If it holds for every $v$, then
\begin{align}
    p_iv=0
    &\Longrightarrow p_ip_jv=0,
    \label{eq:cpgsv21_martingale_overlap_compression_kernel_implication}\\
    p_j(\ker p_i)
    &\subseteq\ker p_i.
    \label{eq:cpgsv21_martingale_overlap_compression_kernel_preservation}
\end{align}
This is stronger than positivity of the smallest nonzero principal angle
between the local ground spaces. A compression proof must therefore derive the
directed estimate
(\ref{eq:cpgsv21_martingale_overlap_eta_norm_compression}); a source-faithful
proof may instead target
(\ref{eq:cpgsv21_martingale_overlap_source_anticommutator}).

The formal development also records a symmetric operator-product route. Assume
\begin{align}
    \lVert p_i(p_jv)\rVert
    &\leq\eta\lVert p_jv\rVert
    \label{eq:cpgsv21_martingale_overlap_operator_product_bound}
\end{align}
for every $v$, equivalently $\lVert p_ip_j\rVert\leq\eta$. By
(\ref{eq:cpgsv21_martingale_overlap_projected_cross_term}),
Cauchy--Schwarz, and the arithmetic--geometric mean inequality,
\begin{align}
    \operatorname{Re}\langle p_iv,p_jv\rangle +\operatorname{Re}\langle p_jv,p_iv\rangle
    &\geq
      -\eta(\lVert p_iv\rVert^2+\lVert p_jv\rVert^2).
    \label{eq:cpgsv21_martingale_overlap_operator_product_anticommutator}
\end{align}
Thus
(\ref{eq:cpgsv21_martingale_overlap_operator_product_bound}) and
(\ref{eq:cpgsv21_martingale_overlap_eta_smallness}) imply
(\ref{eq:cpgsv21_martingale_overlap_eta_gap}). Unlike
(\ref{eq:cpgsv21_martingale_overlap_eta_norm_compression}), this hypothesis
does not force preservation of $\ker p_i$.

The operator-product condition is nevertheless not the source principal-angle
estimate. Here $p_i=h_i$ are excitation projections, whereas
(\ref{eq:cpgsv21_martingale_overlap_ground_projector_defect}) concerns
ground-space projectors. The quantity $\lVert p_ip_j\rVert$ is the cosine of
the smallest principal angle between the excitation ranges. If those ranges
share a nonzero vector, then
\begin{align}
    \lVert p_ip_j\rVert
    &=1.
    \label{eq:cpgsv21_martingale_overlap_excitation_product_unit_norm}
\end{align}
Consequently, the target $\eta<1/(2(L-1))$ cannot follow from a ground-space
intersection estimate without an additional quotient or disjointness
argument. Both norm conditions are conditional formal reductions, not
achievable source-faithful estimates in the generic overlapping case.

The abstract conversion is
\leanid{LinearMap.IsSymmetricProjection.re_inner_anticommutator_ge_neg_of_norm_apply_le}. The cyclic-window
entry points are
\leanid{MPSTensor.parentHamiltonianES_gap_bound_of_cyclic_window_overlap_operator_norm_of_le}
and
\leanid{MPSTensor.parentHamiltonianES_gap_bound_of_cyclic_window_overlap_operator_norm_of_lt};
the public reformulation is
\leanid{MPSTensor.parentHamiltonian_gapped_of_overlap_operator_norm_constant}. For the
blocked range-two theorem, the exact three-site identity and the forward and
backward transport results replace this conditional route. The theorem
\leanid{Submodule.re_inner_anticommutator_starProjection_orthogonal_ge_neg_of_friedrichs_bound}
uses the intersection identity
(\ref{eq:cpgsv21_martingale_overlap_common_ground_space}), equivalently
\leanid{MPSTensor.adjacent_localTermES_eq_zero_iff_mem_groundSpaceES_succ}; it does not
require a small bound on $\lVert p_ip_j\rVert$.

\section{Relation to the blueprint}

The Blueprint item \path{thm:martingale_mps} in
\path{blueprint/src/chapter/ch13_parent_hamiltonian_spectral_gap_martingale.tex}
records the completed blocked range-two theorem. For a primitive tensor with a
positive-definite fixed point, it chooses $p>0$, defines
$B=\operatorname{block}_p(A)$, and proves
\begin{align}
    \operatorname{gap}(H_{N,2}(B))
    &\geq\frac{1}{8},
    &
    N&\geq4.
    \label{eq:cpgsv21_martingale_overlap_blocked_range_two_gap}
\end{align}
The preceding Blueprint item records the exact three-site projector identity
and the transport of the $7/16$ adjacent anticommutator estimate to every
cyclic overlap.

The source-matching entry \path{thm:anticommutator_gap_bound} points to
\leanid{MPSTensor.parentHamiltonianES_gap_bound_of_cyclic_window_overlap_anticommutator}
in \path{TNLean/MPS/ParentHamiltonian/Martingale/Reduction.lean}; its hypothesis
is (\ref{eq:cpgsv21_martingale_overlap_cyclic_anticommutator_target}). The
entries \path{thm:overlap_norm_gap_bound} and
\path{thm:overlap_norm_gap_bound_constant} point to the special and
constant-$\eta$ cyclic-window norm-compression reductions in the same file. The
special reduction uses
\begin{align}
    \eta
    &=
      \left(1-\frac{1}{4L}\right)\frac{1}{2(L-1)},
    \label{eq:cpgsv21_martingale_overlap_special_eta}\\
    1-2\eta(L-1)
    &=\frac{1}{4L}.
    \label{eq:cpgsv21_martingale_overlap_special_eta_gap}
\end{align}
These norm-compression theorems remain conditional on the stronger input. The
finite-row, cyclic-window, non-overlap, anticommutator, norm-compression, and
spectral reductions are formalized, and the blocked range-two MPS
specialization is complete. Comparing the blocked Hamiltonian with an
original-tensor parent interaction in original-site units is tracked under
issue~\#5926.

\section{Conclusion}

The source states the martingale condition with general row-sum coefficients
and invokes principal-angle estimates to prove a positive MPS
parent-Hamiltonian gap~\cite[Section~IV.C]{Cirac2021Matrix}. The formal proof
contains the direct finite-overlap anticommutator reduction for cyclic windows.
For a suitable blocked range-two tensor, the coefficient $7/16$ gives the gap
$1/8$. The stronger conditional norm-compression route assumes
\begin{align}
    \lVert p_ip_jv\rVert
    &\leq\eta\lVert p_iv\rVert,
    &
    2\eta(L-1)&<1.
    \label{eq:cpgsv21_martingale_overlap_conclusion_compression}
\end{align}
The coefficient in
(\ref{eq:cpgsv21_martingale_overlap_special_eta}) is one sufficient choice,
not the only usable martingale coefficient.

The statement of Ref.~\cite{Cirac2021Matrix} is not challenged. The formal
development proves a sufficient explicit estimate for a range-two parent
Hamiltonian after blocking. Recovering the optimized FNW numerator and its
source constants is a separate comparison.

\section{Inverse-Gram reduction and the remaining optimized FNW comparison}

For $n\geq1$, define
\begin{align}
    [n]
    &:=\{0,\ldots,n-1\}.
    \label{eq:cpgsv21_martingale_overlap_finite_index_notation}
\end{align}
Fix an MPS tensor $A=\{A^s\}_{s\in[d]}$ with
$A^s\in\mathbb C^{D\times D}$. For $a,b\in[D]$, let $E_{ab}$ be the matrix
unit with its unique nonzero entry in row $a$ and column $b$. Define the
transfer map and its $L$-fold composition by
\begin{align}
    \mathcal E_A(X)
    &:=\sum_{s\in[d]}A^sX(A^s)^\dagger,
    \label{eq:cpgsv21_martingale_overlap_transfer_map}\\
    \mathcal E_A^L
    &:=\underbrace{\mathcal E_A\circ\cdots\circ\mathcal E_A}
       _{L\text{ factors}}.
    \label{eq:cpgsv21_martingale_overlap_transfer_power}
\end{align}
For a linear map
$T\colon\mathbb C^{D\times D}\to\mathbb C^{D\times D}$, use the rectangular
Choi convention
\begin{align}
    J(T)_{(i,a),(j,b)}
    &:=(T(E_{ij}))_{a,b}.
    \label{eq:cpgsv21_martingale_overlap_rectangular_choi}
\end{align}

Define the open-chain boundary map by
\begin{align}
    \Gamma_L(X)(\sigma)
    &:=\tr(A^\sigma X).
    \label{eq:cpgsv21_martingale_overlap_boundary_map}
\end{align}
Let $U\colon\mathbb C^{D\times D}\to\ell^2([D]\times[D])$ be Frobenius
vectorization, normalized by
\begin{align}
    U(E_{ab})
    &=\mathbf e_{(b,a)}.
    \label{eq:cpgsv21_martingale_overlap_frobenius_vectorization}
\end{align}
After the canonical physical $\ell^2$ identification, let
$\Gamma_L^{\mathrm{ES}}$ denote the transported boundary map and set
\begin{align}
    \mathcal K_L
    &:=
      (\Gamma_L^{\mathrm{ES}})^\dagger\Gamma_L^{\mathrm{ES}}.
    \label{eq:cpgsv21_martingale_overlap_ground_space_gram}
\end{align}
The finite-dimensional realization, range, and Gram operator are
\leanid{MPSTensor.groundSpaceMapES},
\leanid{MPSTensor.range_groundSpaceMapES}, and
\leanid{MPSTensor.groundSpaceGram}. In the chosen coordinate order,
\begin{align}
    \left\langle
      \mathbf e_{(b,a)},\mathcal K_L\mathbf e_{(e,c)}
    \right\rangle_{\ell^2}
    &=J(\mathcal E_A^L)_{(c,e),(a,b)},
    \label{eq:cpgsv21_martingale_overlap_gram_choi_entry}\\
    J(\mathcal E_A^L)_{(c,e),(a,b)}
    &=(\mathcal E_A^L(E_{ca}))_{e,b}.
    \label{eq:cpgsv21_martingale_overlap_gram_transfer_entry}
\end{align}
This is
\leanid{MPSTensor.inner_single_groundSpaceGram_single_eq_rectangularChoi}.

For a matrix endomorphism $\mathcal L$, define its Euclidean reshuffling by
\begin{align}
    \mathscr R(\mathcal L)_{(b,a),(e,c)}
    &:=
      J(\mathcal L)_{(c,e),(a,b)}.
    \label{eq:cpgsv21_martingale_overlap_gram_reshuffling}
\end{align}
The corresponding algebraic matrix is
\leanid{Matrix.gramReshuffleMatrix}, and the operator is
\leanid{Matrix.gramReshuffle}. The theorem
\leanid{Matrix.gramReshuffleMatrix_norm_sq_eq_rectangularChoi_norm_sq} proves
that reshuffling preserves the entrywise Frobenius squared sum.

For every finite set $I$ and $B\in\mathbb C^{I\times I}$,
\leanid{Matrix.toEuclideanCLM_norm_sq_le_sum_norm_sq} gives
\begin{align}
    \lVert B\rVert_{\ell^2(I)\to\ell^2(I)}^2
    &\leq\sum_{i,j\in I}\lvert B_{ij}\rvert^2.
    \label{eq:cpgsv21_martingale_overlap_matrix_operator_frobenius_bound}
\end{align}
Together with reshuffling invariance of the entrywise sum, this gives
\leanid{Matrix.gramReshuffle_norm_sq_le_rectangularChoi_norm_sq}:
\begin{align}
    \lVert\mathscr R(\mathcal L)\rVert_{
      \ell^2(I\times I)\to\ell^2(I\times I)}^2
    &\leq
      \sum_{u,v\in I\times I}\lvert J(\mathcal L)_{u,v}\rvert^2.
    \label{eq:cpgsv21_martingale_overlap_reshuffled_choi_bound}
\end{align}

Equip $\mathbb C^{I\times I}$ with its $L^2$ operator norm and denote the
induced norm of $\mathcal L$ by
$\lVert\mathcal L\rVert_{\mathrm{op},L^2}$. The results
\leanid{Matrix.sum_column_norm_sq_le_norm_sq},
\leanid{Matrix.sum_entry_norm_sq_le_card_mul_norm_sq}, and
\leanid{Matrix.l2_opNorm_single_one} imply
\leanid{Matrix.gramReshuffle_norm_sq_le_card_cube_mul_opNorm_sq}:
\begin{align}
    \lVert\mathscr R(\mathcal L)\rVert_{
      \ell^2(I\times I)\to\ell^2(I\times I)}^2
    &\leq
      |I|^3\lVert\mathcal L\rVert_{\mathrm{op},L^2}^2.
    \label{eq:cpgsv21_martingale_overlap_reshuffling_dimension_bound}
\end{align}
For $I=\Fin{D}$, this gives the norm factor $D^{3/2}$. It is a safe generic
bound, not a sharp estimate or the coefficient in the FNW weighted-Gram
argument.

Reshuffling does not preserve operator norms or singular values. For the
identity map,
\begin{align}
    \lVert\mathrm{id}\rVert_{\mathrm{op},L^2}
    &=1,
    \label{eq:cpgsv21_martingale_overlap_identity_map_norm}\\
    \lVert\mathscr R(\mathrm{id})\rVert
    &=D
    \qquad(D>0).
    \label{eq:cpgsv21_martingale_overlap_identity_reshuffling_norm}
\end{align}
Thus the $D^{3/2}$ factor is an upper bound rather than an invariance result.

The theorem
\leanid{MPSTensor.groundSpaceGram_eq_gramReshuffle} upgrades the coordinate
formula to
\begin{align}
    \mathcal K_L
    &=\mathscr R(\mathcal E_A^L).
    \label{eq:cpgsv21_martingale_overlap_exact_gram_reshuffling}
\end{align}
The ordering is a Choi reshuffling and cannot be replaced by the generally
false Hilbert--Schmidt expression
$\langle E_{ab},\mathcal E_A^L(E_{ce})\rangle_{\mathrm{HS}}$. The theorem
\leanid{Matrix.gramReshuffle_sub} proves compatibility with subtraction.

Let $\rho\in\mathbb C^{D\times D}$ have nonzero trace and define
\begin{align}
    P_\rho(X)
    &:=
      \frac{\tr(X)}{\tr(\rho)}\rho.
    \label{eq:cpgsv21_martingale_overlap_fixed_point_projection}
\end{align}
Assume that $\mathcal E_A$ is trace-preserving,
$\mathcal E_A(\rho)=\rho$, and
\begin{align}
    \rho_{\mathrm{spec}}(\mathcal E_A-P_\rho)
    &<1.
    \label{eq:cpgsv21_martingale_overlap_complement_spectral_radius}
\end{align}
The theorem \leanid{geometric_bound_of_spectralRadius_lt_one} gives
$C>0$ and $0<r<1$ such that
\begin{align}
    \lVert(\mathcal E_A-P_\rho)^n\rVert
    &\leq Cr^n
    \qquad(n\in\mathbb N).
    \label{eq:cpgsv21_martingale_overlap_geometric_transfer_bound}
\end{align}
The transfer-power decomposition, exact reshuffling identity, subtraction
compatibility, and dimension estimate give
\leanid{MPSTensor.groundSpaceGram_sub_fixedPointProj_norm_sq_le_geometric}:
\begin{align}
    \lVert\mathcal K_n-\mathscr R(P_\rho)\rVert^2
    &\leq D^3(Cr^n)^2
    \qquad(n\geq1).
    \label{eq:cpgsv21_martingale_overlap_gram_geometric_bound}
\end{align}
Consequently,
\leanid{MPSTensor.groundSpaceGram_tendsto_gramReshuffle_fixedPointProj} proves
\begin{align}
    \mathcal K_n
    &\longrightarrow\mathscr R(P_\rho)
    \label{eq:cpgsv21_martingale_overlap_gram_limit}
\end{align}
in operator norm. These results require only nonzero trace, trace preservation,
the fixed-point equation, and complementary spectral radius below one. They do
not assume positivity, injectivity, or uniqueness.

The limiting Gram is generally not the identity. Its action is right
multiplication in the vectorization $U$:
\begin{align}
    \mathscr R(P_\rho)U(X)
    &=
      U\left(\frac{X\rho}{\tr\rho}\right).
    \label{eq:cpgsv21_martingale_overlap_limiting_gram_action}
\end{align}
This is
\leanid{Matrix.gramReshuffle_fixedPointProj_frobeniusEquivEuclidean_apply}.
A positive-definite $\rho$ makes the limiting Gram invertible, so perturbation
must be centered at $\mathscr R(P_\rho)$ rather than at the identity.

For an injective finite-dimensional map $T$, the generic range-projector
identity is
\begin{align}
    P_{\operatorname{ran}T}
    &=T(T^\dagger T)^{-1}T^\dagger.
    \label{eq:cpgsv21_martingale_overlap_injective_range_projector}
\end{align}
It is formalized as
\leanid{ContinuousLinearMap.injectiveRangeProjector_eq_starProjection}. The
estimate
\begin{align}
    \lVert(1-t)^{-1}\rVert
    &\leq\frac{1}{1-a}
    \qquad
    \text{if }\lVert t\rVert\leq a<1
    \label{eq:cpgsv21_martingale_overlap_neumann_inverse_bound}
\end{align}
is \leanid{NormedRing.norm_inverse_one_sub_le}.

Set
\begin{align}
    S_T
    &:=(T^\dagger T)^{-1}.
    \label{eq:cpgsv21_martingale_overlap_inverse_gram_definition}
\end{align}
The identities
\leanid{ContinuousLinearMap.inverseGram_comp_adjoint_comp_self} and
\leanid{ContinuousLinearMap.adjoint_comp_self_comp_inverseGram} show that
$S_T$ is a two-sided inverse of $T^\dagger T$, and
\leanid{ContinuousLinearMap.inverseGram_eq_ringInverse} identifies it with the
ring inverse. Hence
\leanid{ContinuousLinearMap.norm_inverseGram_le_of_norm_one_sub_adjoint_comp_self_le}
proves
\begin{align}
    \lVert S_T\rVert
    &\leq\frac{1}{1-a}
    \qquad
    \text{if }\lVert1-T^\dagger T\rVert\leq a<1.
    \label{eq:cpgsv21_martingale_overlap_inverse_gram_identity_bound}
\end{align}

Assume that $L_0>0$ and $A$ is $L_0$-block injective. The theorems
\leanid{MPSTensor.groundSpaceMapES_injective_of_isNBlkInjective} and
\leanid{MPSTensor.groundSpaceMapES_injective_of_isNBlkInjective_of_le} show
that $\Gamma_L^{\mathrm{ES}}$ is injective for every $L\geq L_0$. Therefore
\leanid{MPSTensor.groundSpaceES_starProjection_eq_groundSpaceMapES_comp_inverseGram_comp_adjoint}
gives
\begin{align}
    P_{\mathcal G_L^{\mathrm{ES}}(A)}
    &=
      \Gamma_L^{\mathrm{ES}}\mathcal K_L^{-1}
      (\Gamma_L^{\mathrm{ES}})^\dagger
    \qquad(L\geq L_0).
    \label{eq:cpgsv21_martingale_overlap_physical_inverse_gram_projector}
\end{align}
This is the projector onto
$\operatorname{ran}\Gamma_L^{\mathrm{ES}}$. It is not
\leanid{fixedPointProj}, which acts on the virtual matrix space.

To state the source estimate, let $l\in\mathbb N$ be the overlap length. Let
$0\leq\lambda<1$ satisfy $|\mu|<\lambda$ for every non-unit eigenvalue $\mu$
of the relevant transfer operator, and let $c\geq0$ be its prefactor. The
source permits $c=k^2$, where $k$ is the bond dimension
\cite[Section~6]{Nachtergaele1996Spectral}. Choose $l$ such that
\begin{align}
    c\lambda^l
    &<1.
    \label{eq:cpgsv21_martingale_overlap_source_transfer_smallness}
\end{align}
These are parameters of the quoted FNW estimate, not estimates derived by the
formal development.

The source-specific comparison of the exact inverse-Gram range projectors on
overlapping physical windows remains open. The following prerequisites and
partial results are complete.

\begin{itemize}
    \item The exact Choi reshuffling is
    \leanid{MPSTensor.groundSpaceGram_eq_gramReshuffle} in
    \path{TNLean/MPS/ParentHamiltonian/GroundSpaceGram.lean}; subtraction
    compatibility is \leanid{Matrix.gramReshuffle_sub}.

    \item The geometric power estimate is
    \leanid{geometric_bound_of_spectralRadius_lt_one}. Under the hypotheses
    stated before
    (\ref{eq:cpgsv21_martingale_overlap_complement_spectral_radius}), the
    quantitative and limiting Gram results are
    \leanid{MPSTensor.groundSpaceGram_sub_fixedPointProj_norm_sq_le_geometric}
    and
    \leanid{MPSTensor.groundSpaceGram_tendsto_gramReshuffle_fixedPointProj}.
    Their limit is $\mathscr R(P_\rho)$, and the squared estimate has factor
    $D^3(Cr^n)^2$.

    \item The quantitative inverse results are
    \leanid{MPSTensor.IsPrimitiveMPS.groundSpaceGram_geometric_inverse_bounds}
    and
    \leanid{MPSTensor.IsPrimitiveMPS.groundSpaceMapES_geometric_inverseGram_bounds}.
    Assume that the bond dimension is nonzero and that the primitive
    fixed-point witness $\rho$ is positive definite. Define
    \begin{align}
        \mathcal K_\infty
        &:=\mathscr R(P_\rho),
        \label{eq:cpgsv21_martingale_overlap_limiting_gram}\\
        \mathcal I_\infty
        &:=\mathcal K_\infty^{-1}.
        \label{eq:cpgsv21_martingale_overlap_limiting_inverse_gram}
    \end{align}
    Then there are $g,c,r>0$, with $r<1$ and
    \begin{align}
        c
        &=\lVert\mathcal I_\infty\rVert g,
        \label{eq:cpgsv21_martingale_overlap_inverse_perturbation_constant}\\
        \lVert\mathcal K_n-\mathcal K_\infty\rVert
        &\leq gr^n
        \qquad(n\geq1).
        \label{eq:cpgsv21_martingale_overlap_gram_norm_geometric_bound}
    \end{align}
    If $cr^n<1$, then $\mathcal K_n$ is invertible and
    \begin{align}
        \lVert\mathcal K_n^{-1}\rVert
        &\leq
          \frac{\lVert\mathcal I_\infty\rVert}{1-cr^n},
        \label{eq:cpgsv21_martingale_overlap_inverse_gram_bound}\\
        \lVert\mathcal K_n^{-1}-\mathcal I_\infty\rVert
        &\leq
          \frac{cr^n}{1-cr^n}\lVert\mathcal I_\infty\rVert.
        \label{eq:cpgsv21_martingale_overlap_inverse_gram_difference}
    \end{align}
    At the same lengths, the physical boundary map is injective, its inverse
    Gram is $\mathcal K_n^{-1}$, and it satisfies the same two bounds. Taking
    square roots in
    (\ref{eq:cpgsv21_martingale_overlap_gram_geometric_bound}) gives
    $g=\sqrt{D^3}\,C$ without further dimension loss.

    \item The theorem
    \leanid{MPSTensor.geometric_smallness_at_c3_lengths} proves
    \begin{align}
        cr^l<1
        &\Longrightarrow cr^{l+1}<1,
        \label{eq:cpgsv21_martingale_overlap_smallness_l_plus_one}\\
        cr^l<1
        &\Longrightarrow cr^{K+l}<1,
        \label{eq:cpgsv21_martingale_overlap_smallness_k_plus_l}\\
        cr^l<1
        &\Longrightarrow cr^{K+l+1}<1
        \qquad(K\in\mathbb N).
        \label{eq:cpgsv21_martingale_overlap_smallness_k_plus_l_plus_one}
    \end{align}
    This supplies the three inverse conditions in the C3 correction telescope,
    but it does not prove contraction of the full physical projector defect or
    the coefficient in Equation~(6.1) of
    Ref.~\cite{Nachtergaele1996Spectral}.

    \item The operator-norm--to--entrywise-Frobenius estimate, its
    reshuffled-Choi corollary, and the dimension conversion are
    \leanid{Matrix.toEuclideanCLM_norm_sq_le_sum_norm_sq},
    \leanid{Matrix.gramReshuffle_norm_sq_le_rectangularChoi_norm_sq}, and
    \leanid{Matrix.gramReshuffle_norm_sq_le_card_cube_mul_opNorm_sq}. For
    $I=\Fin{D}$, the last theorem gives norm factor $D^{3/2}$. This factor is
    not claimed sharp or identified with the FNW coefficient.

    \item The spectator-indexed tail and left boundary maps are defined in
    \path{TNLean/MPS/ParentHamiltonian/SpectatorBoundary.lean}. Their Hilbert
    realizations are \leanid{MPSTensor.tailBoundaryMapES} and
    \leanid{MPSTensor.leftBoundaryMapES} in
    \path{TNLean/MPS/ParentHamiltonian/SpectatorBoundaryGram.lean}. That module
    proves their exact ranges, factorizations through
    $\Gamma_{K+L}^{\mathrm{ES}}$, and injectivity from the local boundary map.
    Their Grams are direct sums of $\mathcal K_L$ on tail fibers and
    $\mathcal K_K$ on left fibers, by
    \leanid{MPSTensor.tailBoundaryMapES_adjoint_comp_self_eq_fiberwise_groundSpaceGram}
    and
    \leanid{MPSTensor.leftBoundaryMapES_adjoint_comp_self_eq_fiberwise_groundSpaceGram}.
    Their inverse Grams are fiberwise by
    \leanid{MPSTensor.inverseGram_tailBoundaryMapES_eq_fiberwise_inverseGram}
    and
    \leanid{MPSTensor.inverseGram_leftBoundaryMapES_eq_fiberwise_inverseGram}.
    These exact unweighted identities include empty spectator types and do not
    assert a direct-sum norm estimate.

    \item The C3$'$ whole-increment algebra and a prefix-uniform estimate are
    complete. For suffix increment $Q$, let
    $\Omega_Q=\Fin{d^Q}$ and define
    \begin{align}
        J_Q^{\mathrm{left}}X
        &:=
          (U(A^jX))_{j\in\Omega_Q}.
        \label{eq:cpgsv21_martingale_overlap_left_virtual_map}
    \end{align}
    In this context, set
    \begin{align}
        \mathcal K_n
        &:=
          (\Gamma_n^{\mathrm{ES}})^\dagger\Gamma_n^{\mathrm{ES}},
        \label{eq:cpgsv21_martingale_overlap_whole_increment_finite_gram}\\
        \mathcal K_\infty
        &:=
          \mathscr R(P_\rho).
        \label{eq:cpgsv21_martingale_overlap_whole_increment_limiting_gram}
    \end{align}
    The declarations
    \leanid{MPSTensor.leftVirtualMapES_norm_map_of_leftCanonical} and
    \leanid{MPSTensor.leftVirtualMapES_norm_le_one_of_leftCanonical} prove
    \begin{align}
        \lVert J_Q^{\mathrm{left}}X\rVert
        &=\lVert X\rVert,
        \label{eq:cpgsv21_martingale_overlap_left_virtual_isometry}\\
        \lVert J_Q^{\mathrm{left}}\rVert
        &\leq1.
        \label{eq:cpgsv21_martingale_overlap_left_virtual_norm}
    \end{align}

    The common-ambient reassociation results are
    \leanid{MPSTensor.physicalReassocES},
    \leanid{MPSTensor.physicalReassocES_comp_groundSpaceMapES},
    \leanid{MPSTensor.reassocTailBoundaryMapES},
    \leanid{MPSTensor.reassocTailBoundaryMapES_comp_tailVirtualMapES}, and
    \leanid{MPSTensor.whole_increment_injectiveRangeProjector_residual_eq}.
    They place the tail window of length $L+Q$, the left window of length
    $K+L$ with $Q$ spectators, and the full window of length $K+L+Q$ in one
    Hilbert space.

    The exact mixed-Gram formulas are
    \leanid{MPSTensor.inner_tailBoundaryMapES_adjoint_leftBoundaryMapES_q},
    \leanid{MPSTensor.inner_reassocTailBoundaryMapES_adjoint_leftBoundaryMapES},
    and
    \leanid{MPSTensor.inner_reassocTailBoundaryMapES_adjoint_leftBoundaryMapES_centered}.
    The overlap factors
    \leanid{MPSTensor.wholeIncrementLeftOverlapFactorES} and
    \leanid{MPSTensor.wholeIncrementRightOverlapFactorES} factor
    \leanid{MPSTensor.wholeIncrementCenteredResidualES} through
    $\mathcal K_L-\mathcal K_\infty$. The corresponding declarations are
    \leanid{MPSTensor.wholeIncrementCenteredResidualES_eq_overlapFactors},
    \leanid{MPSTensor.wholeIncrementCenteredResidualES_norm_le},
    \leanid{MPSTensor.wholeIncrementCenteredResidualES_norm_le_of_leftCanonical},
    and
    \leanid{MPSTensor.IsPrimitiveMPS.wholeIncrementCenteredResidualES_norm_le}.

    The finite-Gram telescope uses
    \leanid{MPSTensor.reassocTailBoundaryMapES_adjoint_comp_self_eq_fiberwise_groundSpaceGram},
    \leanid{MPSTensor.inverseGram_reassocTailBoundaryMapES_eq_fiberwise_inverseGram},
    \leanid{MPSTensor.wholeIncrementLeftFiniteGramCancellation_eq_groundSpaceMapES_adjoint},
    \leanid{MPSTensor.wholeIncrementTailFiniteGramCancellation_eq_groundSpaceMapES_adjoint},
    and
    \leanid{MPSTensor.wholeIncrement_limitingOverlapProduct_eq}.
    Its correction maps are
    \leanid{MPSTensor.wholeIncrementTailFiniteGramCorrectionES},
    \leanid{MPSTensor.wholeIncrementFullInverseGramCorrectionES}, and
    \leanid{MPSTensor.wholeIncrementLeftFiniteGramCorrectionES}. The exact
    telescopes are
    \leanid{MPSTensor.wholeIncrement_virtualResidual_eq_centered_sub_gramErrors}
    and
    \leanid{MPSTensor.wholeIncrement_injectiveRangeProjector_residual_eq_centered_sub_corrections},
    the latter using
    \leanid{MPSTensor.wholeIncrementCenteredProjectorResidualES}. The
    pseudoinverse and centered-sandwich estimates are
    \leanid{MPSTensor.norm_reassocTailBoundaryMapES_comp_inverseGram_le_sqrt}
    and
    \leanid{MPSTensor.wholeIncrementCenteredProjectorResidualES_norm_le}. The
    correction estimates are
    \leanid{MPSTensor.wholeIncrementTailFiniteGramCorrectionES_norm_le},
    \leanid{MPSTensor.wholeIncrementFullInverseGramCorrectionES_norm_le}, and
    \leanid{MPSTensor.wholeIncrementLeftFiniteGramCorrectionES_norm_le}.

    The assembled theorem is
    \leanid{MPSTensor.IsPrimitiveMPS.wholeIncrement_groundProjection_defect_le_geometric}.
    For a primitive tensor with positive-definite fixed point and nonzero bond
    dimension, it gives $C,c,r>0$, with $r<1$, uniformly in the prefix length
    $K$ and suffix length $Q$. If $L\geq1$ and $cr^L<1$, define
    \begin{align}
        P^{\mathrm{tail}}_{K;L,Q}
        &:=
          P_{\operatorname{ran}
          (\widetilde\Gamma^{\mathrm{tail}}_{K;L,Q})},
        \label{eq:cpgsv21_martingale_overlap_whole_increment_tail_projector}\\
        P^{\mathrm{left}}_{K+L,Q}
        &:=
          P_{\operatorname{ran}
          (\Gamma^{\mathrm{left}}_{K+L,Q})}.
        \label{eq:cpgsv21_martingale_overlap_whole_increment_left_projector}
    \end{align}
    Then
    \begin{align}
        \left\lVert
          P^{\mathrm{tail}}_{K;L,Q}P^{\mathrm{left}}_{K+L,Q}
          -P_{\mathcal G_{K+L+Q}^{\mathrm{ES}}(A)}
        \right\rVert
        &\leq\frac{Cr^L}{1-cr^L}.
        \label{eq:cpgsv21_martingale_overlap_whole_increment_defect}
    \end{align}
    This reconstructed estimate is not attributed to FNW and does not claim
    the optimized numerator in Equation~(6.1) of
    Ref.~\cite{Nachtergaele1996Spectral}.

    The theorem
    \leanid{MPSTensor.IsPrimitiveMPS.exists_uniform_wholeIncrement_defect_le_seven_sixteenths}
    chooses $p>0$, measured in input-tensor sites, such that $A$ is
    $p$-block injective and the defect with $L=Q=p$ is at most $7/16$ for
    every $K$. Its specialization
    \leanid{MPSTensor.IsPrimitiveMPS.exists_threeBlock_wholeIncrement_defect_le_seven_sixteenths}
    takes $K=p$. Thus the prefix-uniform estimate and the three-block
    coefficient $7/16$ are proved in original-site units.

    \item The mixed Gram of the C3 tail map at $(K,l+1)$ and the left map at
    $(K+l,1)$ is computed by
    \leanid{MPSTensor.inner_tailBoundaryMapES_adjoint_leftBoundaryMapES}.
    Write these maps as $\Gamma^{\mathrm{tail}}_{K,l+1}$ and
    $\Gamma^{\mathrm{left}}_{K+l,1}$. For boundary-family vectors $x,y$ with
    matrix fibers $X_u,Z_j$,
    \begin{align}
        \left\langle
          x,
          (\Gamma^{\mathrm{tail}}_{K,l+1})^\dagger
          \Gamma^{\mathrm{left}}_{K+l,1}y
        \right\rangle
        &=
          \sum_{u\in\Fin{d^K}}\sum_{j\in\Fin{d}}
          \left\langle
            U(A^jX_u),\mathcal K_lU(Z_jA^u)
          \right\rangle.
        \label{eq:cpgsv21_martingale_overlap_c3_mixed_gram}
    \end{align}
    This reconstructed algebraic identity follows from the boundary-map
    definitions, word factorization, and trace cyclicity. It is not attributed
    to FNW or Nachtergaele. The sourced target is the projector defect in
    Equation~(2.4), with the rational estimate stated after that equation and
    in Equation~(6.1), of
    Ref.~\cite{Nachtergaele1996Spectral}. No mixed-Gram norm estimate follows
    from (\ref{eq:cpgsv21_martingale_overlap_c3_mixed_gram}) alone.

    \item For injective maps with common factorizations
    \begin{align}
        TJ_1
        &=C,
        \label{eq:cpgsv21_martingale_overlap_common_factorization_tail}\\
        LJ_2
        &=C,
        \label{eq:cpgsv21_martingale_overlap_common_factorization_left}
    \end{align}
    the theorem
    \leanid{ContinuousLinearMap.injectiveRangeProjector_comp_sub_of_factorizations}
    factors the projector residual through
    \begin{align}
        T^\dagger L
        -(T^\dagger T)J_1S_CJ_2^\dagger(L^\dagger L).
        \label{eq:cpgsv21_martingale_overlap_projector_residual_virtual_factor}
    \end{align}
    The specialization
    \leanid{MPSTensor.c3_injectiveRangeProjector_residual_eq} uses the tail
    map at $(K,L_0+1)$, the left map at $(K+L_0,1)$, the common map
    $\Gamma_{K+L_0+1}^{\mathrm{ES}}$, and the corresponding virtual maps. Its
    three inverse-Gram projectors have the open-chain tail, left, and full
    ground-space ranges. All injectivity assumptions are explicit. This is an
    algebraic identity; a source-specific mixed-Gram estimate is still needed.

    \item For injective $T$ and $S_T=(T^\dagger T)^{-1}$,
    \leanid{ContinuousLinearMap.norm_T_comp_inverseGram_eq_sqrt} and
    \leanid{ContinuousLinearMap.norm_inverseGram_comp_adjoint_eq_sqrt} prove
    \begin{align}
        \lVert TS_T\rVert
        &=\sqrt{\lVert S_T\rVert},
        \label{eq:cpgsv21_martingale_overlap_left_pseudoinverse_norm}\\
        \lVert S_TT^\dagger\rVert
        &=\sqrt{\lVert S_T\rVert}.
        \label{eq:cpgsv21_martingale_overlap_right_pseudoinverse_norm}
    \end{align}

    \item For a primitive MPS tensor with nonzero virtual dimension, let
    $J_K^{\mathrm{tail}}$ be its length-$K$ tail virtual map. The theorem
    \leanid{MPSTensor.IsPrimitiveMPS.tailVirtualMapES_norm_uniform} gives
    $C>0$, independent of $K$, such that
    \begin{align}
        \lVert J_K^{\mathrm{tail}}\rVert
        &\leq C
        \qquad(K\in\mathbb N).
        \label{eq:cpgsv21_martingale_overlap_uniform_tail_virtual_norm}
    \end{align}
    It does not assume positive-definiteness of $\rho$. The proof first obtains
    a uniform fourth-power estimate and then takes fourth roots. It uses the
    exact tail adjoint--self estimate, the fixed-point/complementary-power
    decomposition, the geometric power estimate, and the conversion from the
    $L^2$ operator norm to entrywise mass. This is not an FNW or Nachtergaele
    contraction estimate.

    \item Let $\Omega_K$ be the set of length-$K$ physical words and let
    $\mathcal B_K=\bigoplus_{u\in\Omega_K}\MN{D}$. Define
    \begin{align}
        J_1^{\mathrm{left}}(X)
        &:=(A^jX)_{j\in\Omega_1},
        \label{eq:cpgsv21_martingale_overlap_one_site_left_virtual_map}\\
        J_K^{\mathrm{tail}}(X)
        &:=(XA^u)_{u\in\Omega_K}.
        \label{eq:cpgsv21_martingale_overlap_k_site_tail_virtual_map}
    \end{align}
    The factors \leanid{MPSTensor.c3LeftOverlapFactorES} and
    \leanid{MPSTensor.c3RightOverlapFactorES} satisfy
    \begin{align}
        (F_KX)_{u,j}
        &=A^jX_u,
        \label{eq:cpgsv21_martingale_overlap_c3_left_overlap_factor}\\
        (G_KZ)_{u,j}
        &=Z_jA^u,
        \label{eq:cpgsv21_martingale_overlap_c3_right_overlap_factor}
    \end{align}
    by
    \leanid{MPSTensor.boundaryFamilyEquiv_c3LeftOverlapFactorES_apply} and
    \leanid{MPSTensor.boundaryFamilyEquiv_c3RightOverlapFactorES_apply}.
    These formulas allow finite, possibly empty, configuration sets. The
    theorems \leanid{MPSTensor.c3LeftOverlapFactorES_norm_le} and
    \leanid{MPSTensor.c3RightOverlapFactorES_norm_le} give
    \begin{align}
        \lVert F_K\rVert
        &\leq\lVert J_1^{\mathrm{left}}\rVert,
        \label{eq:cpgsv21_martingale_overlap_c3_left_factor_norm}\\
        \lVert G_K\rVert
        &\leq\lVert J_K^{\mathrm{tail}}\rVert.
        \label{eq:cpgsv21_martingale_overlap_c3_right_factor_norm}
    \end{align}
    No prefix-cardinality factor appears.

    If $D\geq1$ and $\rho$ is positive definite, set
    $\mathcal K_\infty=\mathscr R(P_\rho)$. The centered residual
    \leanid{MPSTensor.c3CenteredVirtualResidualES} factors as
    \begin{align}
        R_{K,l}
        &=
          F_K^\dagger
          (\mathcal K_l-\mathcal K_\infty)^{
             \oplus(\Omega_K\times\Omega_1)}
          G_K.
        \label{eq:cpgsv21_martingale_overlap_c3_centered_factorization}
    \end{align}
    This is
    \leanid{MPSTensor.c3_centeredVirtualResidualES_eq_overlapFactors}.
    The theorem \leanid{MPSTensor.c3_centeredVirtualResidualES_norm_le} gives
    \begin{align}
        \lVert R_{K,l}\rVert
        &\leq
          \lVert J_1^{\mathrm{left}}\rVert
          \lVert\mathcal K_l-\mathcal K_\infty\rVert
          \lVert J_K^{\mathrm{tail}}\rVert.
        \label{eq:cpgsv21_martingale_overlap_c3_centered_norm}
    \end{align}
    If $A$ is primitive with positive-definite fixed-point witness $\rho$, then
    \leanid{MPSTensor.IsPrimitiveMPS.c3_centeredVirtualResidualES_norm_sq_le_geometric}
    gives $C>0$ and $0<r<1$, independent of $K$, such that
    \begin{align}
        \lVert R_{K,l}\rVert^2
        &\leq D^3(Cr^l)^2
        \qquad(K\in\mathbb N,\ l\geq1).
        \label{eq:cpgsv21_martingale_overlap_c3_centered_geometric_bound}
    \end{align}
    These reconstructed identities control only the centered term. They prove
    neither the full defect nor the rational coefficient in Equation~(6.1) of
    Ref.~\cite{Nachtergaele1996Spectral}.

    \item Whenever $\mathcal K_n$ is invertible, set
    $S_n=\mathcal K_n^{-1}$. For positive-definite $\rho$, set
    $\mathcal K_\infty=\mathscr R(P_\rho)$ and
    $S_\infty=\mathcal K_\infty^{-1}$. The theorem
    \leanid{MPSTensor.c3_virtualResidual_eq_centered_sub_gramErrors} telescopes
    the exact C3 residual into the centered residual and three corrections
    containing
    \begin{align}
        \mathcal K_{l+1}-\mathcal K_\infty,
        \label{eq:cpgsv21_martingale_overlap_c3_tail_gram_error}\\
        S_{K+l+1}-S_\infty,
        \label{eq:cpgsv21_martingale_overlap_c3_full_inverse_error}\\
        \mathcal K_{K+l}-\mathcal K_\infty.
        \label{eq:cpgsv21_martingale_overlap_c3_left_gram_error}
    \end{align}
    The theorem
    \leanid{MPSTensor.c3_injectiveRangeProjector_residual_eq_centered_sub_gramErrors}
    places this identity between the local pseudoinverse factors.

    Under nonzero virtual dimension and positive-definite $\rho$,
    \leanid{MPSTensor.inner_limitingMixedGramIntertwining} proves the limiting
    intertwining identity, and
    \leanid{MPSTensor.inner_c3_centeredVirtualResidual_eq_gramError} identifies
    the centered residual pairing with the length-$l$ Gram error. The
    cancellation uses the nonidentity limiting metric and its inverse; it does
    not require a transfer fixed-point equation.

    \item The physical corrections are
    \leanid{MPSTensor.c3TailFiniteGramCorrectionES},
    \leanid{MPSTensor.c3FullInverseGramCorrectionES}, and
    \leanid{MPSTensor.c3LeftFiniteGramCorrectionES}. Write
    \begin{align}
        T
        &:=\Gamma^{\mathrm{tail}}_{K,l+1},
        \label{eq:cpgsv21_martingale_overlap_c3_tail_map}\\
        L
        &:=\Gamma^{\mathrm{left}}_{K+l,1},
        \label{eq:cpgsv21_martingale_overlap_c3_left_map}\\
        C
        &:=\Gamma_{K+l+1}^{\mathrm{ES}}.
        \label{eq:cpgsv21_martingale_overlap_c3_full_map}
    \end{align}
    The exact factorizations
    \leanid{MPSTensor.c3TailFiniteGramCorrectionES_eq_factored},
    \leanid{MPSTensor.c3FullInverseGramCorrectionES_eq_factored}, and
    \leanid{MPSTensor.c3LeftFiniteGramCorrectionES_eq_factored} are
    \begin{align}
        Q_1
        &=
          TS_{l+1}^{\oplus\Omega_K}
          (\mathcal K_{l+1}-\mathcal K_\infty)^{\oplus\Omega_K}
          J_K^{\mathrm{tail}}S_{K+l+1}C^\dagger,
        \label{eq:cpgsv21_martingale_overlap_c3_first_correction}\\
        Q_2
        &=
          TS_{l+1}^{\oplus\Omega_K}
          \mathcal K_\infty^{\oplus\Omega_K}
          J_K^{\mathrm{tail}}S_\infty
          (\mathcal K_\infty-\mathcal K_{K+l+1})
          S_{K+l+1}C^\dagger,
        \label{eq:cpgsv21_martingale_overlap_c3_second_correction}\\
        Q_3
        &=
          TS_{l+1}^{\oplus\Omega_K}
          \mathcal K_\infty^{\oplus\Omega_K}
          J_K^{\mathrm{tail}}S_\infty
          (J_1^{\mathrm{left}})^\dagger
          (\mathcal K_{K+l}-\mathcal K_\infty)^{\oplus\Omega_1}
          S_{K+l}^{\oplus\Omega_1}L^\dagger.
        \label{eq:cpgsv21_martingale_overlap_c3_third_correction}
    \end{align}
    The second line uses
    \leanid{MPSTensor.inverseGram_sub_limitingInverse_eq_resolvent}:
    \begin{align}
        S_{K+l+1}-S_\infty
        &=
          S_\infty
          (\mathcal K_\infty-\mathcal K_{K+l+1})
          S_{K+l+1}.
        \label{eq:cpgsv21_martingale_overlap_inverse_resolvent_identity}
    \end{align}
    The fiberwise subtraction theorem is
    \leanid{MPSTensor.boundaryFiberwiseMap_sub}, and
    \leanid{MPSTensor.groundSpaceGram_isUnit_of_groundSpaceMapES_injective}
    derives Gram invertibility from injectivity.

    The spectator pseudoinverse estimates
    \leanid{MPSTensor.norm_tailBoundaryMapES_comp_inverseGram_le_sqrt} and
    \leanid{MPSTensor.norm_inverseGram_comp_leftBoundaryMapES_adjoint_le_sqrt},
    together with the fiberwise norm estimate, yield
    \leanid{MPSTensor.c3TailFiniteGramCorrectionES_norm_le},
    \leanid{MPSTensor.c3FullInverseGramCorrectionES_norm_le}, and
    \leanid{MPSTensor.c3LeftFiniteGramCorrectionES_norm_le}:
    \begin{align}
        \lVert Q_1\rVert
        &\leq
          \sqrt{\lVert S_{l+1}\rVert}
          \lVert\mathcal K_{l+1}-\mathcal K_\infty\rVert
          \lVert J_K^{\mathrm{tail}}\rVert
          \sqrt{\lVert S_{K+l+1}\rVert},
        \label{eq:cpgsv21_martingale_overlap_c3_first_correction_bound}\\
        \lVert Q_2\rVert
        &\leq
          \sqrt{\lVert S_{l+1}\rVert}
          \lVert\mathcal K_\infty\rVert
          \lVert J_K^{\mathrm{tail}}\rVert
          \lVert S_\infty\rVert
          \lVert\mathcal K_{K+l+1}-\mathcal K_\infty\rVert
          \sqrt{\lVert S_{K+l+1}\rVert},
        \label{eq:cpgsv21_martingale_overlap_c3_second_correction_bound}\\
        \lVert Q_3\rVert
        &\leq
          \sqrt{\lVert S_{l+1}\rVert}
          \lVert\mathcal K_\infty\rVert
          \lVert J_K^{\mathrm{tail}}\rVert
          \lVert S_\infty\rVert
          \lVert J_1^{\mathrm{left}}\rVert
          \lVert\mathcal K_{K+l}-\mathcal K_\infty\rVert
          \sqrt{\lVert S_{K+l}\rVert}.
        \label{eq:cpgsv21_martingale_overlap_c3_third_correction_bound}
    \end{align}
    No spectator-cardinality factor occurs. These formulas follow from the
    inverse-Gram representations of the physical projectors. They are not
    formulas stated by
    Refs.~\cite{Fannes1992Finitely,Nachtergaele1996Spectral,Cirac2021Matrix}.
\end{itemize}

Fix $K,L_0\in\mathbb N$ and define the common ambient Hilbert space
\begin{align}
    \mathcal H_{K+L_0+1}
    &:=
      \ell^2(\Fin{d^{K+L_0+1}}).
    \label{eq:cpgsv21_martingale_overlap_c3_common_ambient}
\end{align}
The tail transport theorem at $(K,L_0+1)$ identifies its range with
$\mathcal V_{K,L_0+1}(A)$. The left transport theorem at $(K+L_0,1)$
identifies its range with $\mathcal U_{K+L_0}(A)$; length one is the left
boundary-map convention matching Nachtergaele's C3 condition
\cite[Equation~(2.4)]{Nachtergaele1996Spectral}. The corresponding defect is
\begin{align}
    \left\lVert
      P^{\mathrm{tail}}_{K,L_0+1}(A)
      P^{\mathrm{left}}_{K+L_0}(A)
      -P_{\mathcal G_{K+L_0+1}^{\mathrm{ES}}(A)}
    \right\rVert.
    \label{eq:cpgsv21_martingale_overlap_c3_physical_defect}
\end{align}
Here $\mathcal U_{K+L_0}(A)$ and $\mathcal V_{K,L_0+1}(A)$ are the left and
tail ground spaces represented by
\leanid{MPSTensor.openChainLeftGroundSpaceES} and
\leanid{MPSTensor.openChainTailGroundSpaceES}. Under block injectivity,
$D\geq1$, and $L_0>0$,
\leanid{MPSTensor.openChainLeft_inf_tailGroundSpaceES} proves
\begin{align}
    \mathcal U_{K+L_0}(A)\cap\mathcal V_{K,L_0+1}(A)
    &=
      \mathcal G_{K+L_0+1}^{\mathrm{ES}}(A).
    \label{eq:cpgsv21_martingale_overlap_c3_ground_space_intersection}
\end{align}
The projectors in
(\ref{eq:cpgsv21_martingale_overlap_c3_physical_defect}) are the orthogonal
projectors onto these spaces.

The common-ambient range identifications are
\leanid{MPSTensor.tailBoundaryMap_range_map_symm_eq_openChainTailGroundSpaceES},
specialized to $(K,L_0+1)$, and
\leanid{MPSTensor.leftBoundaryMap_range_map_symm_eq_openChainLeftGroundSpaceES},
specialized to $(K+L_0,1)$.

The formal estimate controls deviation from the nonidentity limiting Gram:
\begin{align}
    \lVert\mathcal K_n-\mathscr R(P_\rho)\rVert^2
    &\leq D^3(Cr^n)^2.
    \label{eq:cpgsv21_martingale_overlap_repeated_gram_control}
\end{align}
It does not control $\lVert1-\mathcal K_n\rVert$ in unweighted coordinates.
Under positive-definiteness of $\rho$, the results
\leanid{MPSTensor.IsPrimitiveMPS.groundSpaceGram_geometric_inverse_bounds} and
\leanid{MPSTensor.IsPrimitiveMPS.groundSpaceMapES_geometric_inverseGram_bounds}
perturb around $\mathscr R(P_\rho)$ and give denominator $1-cr^n$ whenever
$cr^n<1$. The theorem
\leanid{MPSTensor.geometric_smallness_at_c3_lengths} supplies the three
required lengths. The square-root conversion retains the factor
$\sqrt{D^3}\,C$ without further dimension loss. Moreover,
\leanid{MPSTensor.IsPrimitiveMPS.groundSpaceGram_ringInverse_tendsto} proves
\begin{align}
    \mathcal K_n^{-1}
    &\longrightarrow\mathscr R(P_\rho)^{-1}.
    \label{eq:cpgsv21_martingale_overlap_inverse_gram_convergence}
\end{align}

For the uniform assembly, set
\begin{align}
    \mathcal K_\infty
    &:=\mathscr R(P_\rho),
    \label{eq:cpgsv21_martingale_overlap_uniform_limiting_gram}\\
    \mathcal I_\infty
    &:=\mathcal K_\infty^{-1}.
    \label{eq:cpgsv21_martingale_overlap_uniform_limiting_inverse}
\end{align}
Choose $g,c,r>0$ from
\leanid{MPSTensor.IsPrimitiveMPS.groundSpaceMapES_geometric_inverseGram_bounds}
with
\begin{align}
    c
    &=\lVert\mathcal I_\infty\rVert g,
    &
    0&<r<1,
    \label{eq:cpgsv21_martingale_overlap_uniform_geometric_constants}
\end{align}
and choose $T>0$ from
\leanid{MPSTensor.IsPrimitiveMPS.tailVirtualMapES_norm_uniform}. Define
\begin{align}
    J
    &:=\lVert J_1^{\mathrm{left}}\rVert,
    \label{eq:cpgsv21_martingale_overlap_uniform_left_constant}\\
    M
    &:=
      \lVert\mathcal K_\infty\rVert
      \lVert\mathcal I_\infty\rVert,
    \label{eq:cpgsv21_martingale_overlap_uniform_condition_number}\\
    C_*
    &:=
      4\lVert\mathcal I_\infty\rVert Tg(J+1)(M+1).
    \label{eq:cpgsv21_martingale_overlap_uniform_defect_constant}
\end{align}
Then
\leanid{MPSTensor.IsPrimitiveMPS.openChain_groundProjection_defect_le_geometric}
proves, for every $K$ and every $l\geq1$ with $cr^l<1$,
\begin{align}
    \left\lVert
      P^{\mathrm{tail}}_{K,l+1}
      P^{\mathrm{left}}_{K+l}
      -P_{\mathcal G_{K+l+1}^{\mathrm{ES}}(A)}
    \right\rVert
    &\leq\frac{C_*r^l}{1-cr^l}.
    \label{eq:cpgsv21_martingale_overlap_open_chain_uniform_defect}
\end{align}
The proof uses
\leanid{MPSTensor.c3_injectiveRangeProjector_residual_eq_centered_sub_corrections},
\leanid{MPSTensor.c3CenteredProjectorResidualES_norm_le}, and
\leanid{MPSTensor.c3_centeredVirtualResidualES_norm_le}. Its other inputs are
the uniform tail estimate, the pointwise inverse-Gram estimates, the
three-length smallness theorem, and the correction estimates
\leanid{MPSTensor.c3TailFiniteGramCorrectionES_norm_le},
\leanid{MPSTensor.c3FullInverseGramCorrectionES_norm_le}, and
\leanid{MPSTensor.c3LeftFiniteGramCorrectionES_norm_le}. Exact range
identities identify the tail, left, and full projectors. No physical projector
is identified with a transfer fixed-point projector.

This reconstructed coefficient suffices for Nachtergaele's C3 condition. The
theorem
\leanid{MPSTensor.IsPrimitiveMPS.exists_openChain_groundProjection_defect_lt_c3_threshold}
uses
\begin{align}
    \sqrt{n+1}\,r^n
    &\longrightarrow0
    \label{eq:cpgsv21_martingale_overlap_c3_decay_threshold}
\end{align}
to choose $l>1$ above a positive block-injectivity length such that
\begin{align}
    cr^l
    &<1,
    \label{eq:cpgsv21_martingale_overlap_c3_chosen_smallness}\\
    \frac{C_*r^l}{1-cr^l}
    &<\frac{1}{\sqrt{l+1}}.
    \label{eq:cpgsv21_martingale_overlap_c3_chosen_threshold}
\end{align}
The same $l$ works for every spectator length $K$. The theorem
\leanid{MPSTensor.IsPrimitiveMPS.exists_re_inner_openChain_anticommutator_ge_c3_threshold}
then gives the open-chain anticommutator estimate through two-projection
geometry. The length $l$ is measured in sites of the input tensor; if that
tensor is already blocked, one such site is one blocked-chain filtration step.

The reconstructed coefficient is not the optimized FNW expression
\begin{align}
    c\lambda^l\frac{1+c\lambda^l}{1-c\lambda^l},
    \qquad
    c\lambda^l&<1,
    \label{eq:cpgsv21_martingale_overlap_optimized_fnw_expression}
\end{align}
quoted as \path{boundAm} in Section~6 of
Ref.~\cite{Nachtergaele1996Spectral}. Recovering its numerator and source
constants requires the weighted Gram estimates in Lemmas~5.2, 5.3, and 6.2 of
Ref.~\cite{Fannes1992Finitely}. The reconstructed constant is not identified
with a source dimension-dependent value; the available unweighted conversion
retains the proved $D^{3/2}$ norm factor inside $g$.

\section{Verdict}

\textbf{Verdict: scope restriction.} The exact whole-increment algebra, the
reconstructed prefix-uniform geometric C3$'$-type estimate, the three-block
coefficient $7/16$, and the resulting cyclic range-two gap $1/8$ are proved
for a suitable blocked tensor. Before blocking, lengths are measured in sites
of the input tensor; the final chain length counts blocked sites. These bounds
come from the formal inverse-Gram reconstruction and are not attributed to
FNW. The optimized FNW numerator and its source constants remain open.

The blocked-to-original comparison is proved for periodic lengths divisible by
the chosen block length. Under the canonical configuration isometry, each
blocked range-two term equals the original range-$2p$ term beginning at a block
boundary. The conjugated blocked Hamiltonian is therefore the sparse sum over
aligned starts. The omitted original local projections form a non-negative
quadratic-form remainder. Block injectivity identifies the blocked and original
kernels with the same periodic MPS line under this isometry. Hence
\leanid{MPSTensor.IsPrimitiveMPS.exists_parentHamiltonianES_gap_eighth_mul}
proves
\begin{align}
    \operatorname{gap}(H_{Np,2p}(A))
    &\geq\frac{1}{8}
    \qquad(N\geq4).
    \label{eq:cpgsv21_martingale_overlap_original_aligned_length_gap}
\end{align}

This finite cyclic sparse-sum comparison is a TNLean reconstruction.
Nachtergaele formulates compatible open intervals in
Equations~(3.12)--(3.16), and the proof of Theorem~2.1(ii) uses the rescaled
sequence $\widetilde\Lambda_n=\Lambda_{ln}$
\cite{Nachtergaele1996Spectral}. That source does not state the finite periodic
operator comparison used here. The remaining case for the unrestricted normal
MPS theorem consists of chain lengths
\begin{align}
    Np+r,
    \qquad
    0&<r<p.
    \label{eq:cpgsv21_martingale_overlap_remainder_lengths}
\end{align}
It requires a boundary or remainder comparison that preserves the common
ground-space complement and a uniform positive constant. No all-length claim
is made.

\bibliographystyle{alpha}
\bibliography{references}

\end{document}
